% Source: https://artofproblemsolving.com/community/c3h1703437_help_me_on_this_problem_its_so_weird
% Stable identity: aops-topic-1703437-q01
Starting at $P_0$ and stepping $m$ places each time, the path visits $P_0,P_m,P_{2m},\dots$ (indices modulo $36$). It returns to $P_0$ for the first time after $\dfrac{36}{\gcd(m,36)}$ steps, so it uses all $36$ points exactly when $\gcd(m,36)=1$.

The values $1\le m\le35$ with $\gcd(m,36)=1$ are the numbers not divisible by $2$ or $3$:
\[
1,5,7,11,13,17,19,23,25,29,31,35,
\]
which is $\varphi(36)=36\cdot\dfrac{1}{2}\cdot\dfrac{2}{3}=12$ values. Stepping $m$ places clockwise draws the same chords as stepping $36-m$ places (that is, $m$ places counterclockwise), so the values pair up as $\{1,35\},\{5,31\},\{7,29\},\{11,25\},\{13,23\},\{17,19\}$, and different pairs give chords of different lengths, hence different figures.


\begin{asy}
usepackage("amsmath");
size(5cm);
defaultpen(fontsize(10pt));
draw(unitcircle,gray+linewidth(0.5));
for(int i=0;i<36;++i){draw(dir(90-10*i)--dir(90-10*(i+13)),blue+linewidth(0.5));dot(dir(90-10*i),linewidth(2.5));}
label("$P_0$",dir(90),N);label("$P_{13}$",dir(-40),SE);
label("\mbox{}",(min(currentpicture,user=true)+max(currentpicture,user=true))/2,black); // keeps the text after the figure black when LaTeX typesets the labels inline
\end{asy}

(The figure shows the star polygon for $m=13$: every point is used, and the path closes after $36$ chords.)

So there are $\boxed{6}$ distinct star polygons. This agrees with the answer book; the guess $7$ is too large. (The pair $\{1,35\}$ gives the regular $36$-gon itself. It uses all $36$ points, so it counts under the stated definition; if it were excluded the count would be $5$.)
